% Research CV for graduate applications (thesis MSc, or PhD). Build: tectonic Sadeeq_Adeyemi_CV_Research.tex
% Content rules (2026-10-14 research, U of T / JHU / UCalgary guidance): education first with grades, thesis and
% supervisor; research entries name the lab and the advisor and say what I did myself; outputs listed in full;
% referees on a separate final page with full contact details ("available on request" is not used in Canada).
% To tailor for a supervisor: put the research interest closest to their lab first.
% Advisor, GPAs and referees left out by choice (2026-10-14); give them when an application asks.
\documentclass[10pt,letterpaper]{article}
% Shared look for both CVs. Single column (applicant tracking systems read it in order), serif body, small
% blue section heads, dates in grey on the right. Built with tectonic (XeTeX). No tables, no text boxes, no icons.
\usepackage[letterpaper,top=0.5in,bottom=0.5in,left=0.6in,right=0.6in]{geometry}
\usepackage{fontspec}
% Merriweather body, Open Sans headings (the pairing of the example layout). Tested: both pypdf and pdfminer read
% every word and "AWS" correctly; Source Serif, Latin Modern, TeX Gyre and Libertinus did not.
\usepackage{merriweather}
\usepackage{opensans}
\newcommand{\headfont}{\sffamily}
\usepackage{enumitem,titlesec,xcolor}
\usepackage[hidelinks]{hyperref}
\pagestyle{empty}
\setlength{\parindent}{0pt}\setlength{\parskip}{0pt}
% Left-aligned, no hyphenation: justified lines squeeze the spaces between words, and some applicant tracking
% systems then read words run together; hyphens split keywords.
\AtBeginDocument{\raggedright}\hyphenpenalty=10000\exhyphenpenalty=10000\sloppy
\definecolor{ink}{HTML}{1A1F27}\definecolor{accent}{HTML}{1F5F8B}\definecolor{grey}{HTML}{5F6873}
\color{ink}

\newcommand{\BodySize}{9.3}
\titleformat{\section}{\headfont\bfseries\fontsize{10.5}{12}\selectfont\color{accent}\MakeUppercase}{}{0pt}{}
\titlespacing*{\section}{0pt}{9pt}{3pt}
\setlist[itemize]{leftmargin=1.05em,itemsep=1.4pt,topsep=1.5pt,parsep=0pt,partopsep=0pt,label={\color{grey}\textendash}}

% \entry{Title}{Organisation, place}{Dates}
\newcommand{\entry}[3]{\vspace{4pt}{\bfseries #1}\ifx\relax#2\relax\else, #2\fi\hfill{\color{grey}\small #3}\par}
% \skillrow{Label}{Items}
\newcommand{\skillrow}[2]{{\bfseries #1:} #2\par\vspace{1.5pt}}
\newcommand{\sep}{\enspace{\color{grey}|}\enspace}
\newcommand{\head}[2]{%
  {\fontsize{22}{26}\selectfont\bfseries Sadeeq Adeyemi}\par\vspace{2pt}
  {\headfont\fontsize{11.5}{14}\selectfont\color{accent} #1}\par\vspace{5pt}
  {\small #2}\par}

\renewcommand{\BodySize}{9.05}
\begin{document}
\newgeometry{top=0.42in,bottom=0.4in,left=0.55in,right=0.55in}
\fontsize{\BodySize}{12.6}\selectfont

\head{3D perception and multi-object tracking under occlusion}{%
Bangor, Maine \sep \href{mailto:sideeqadeyemi@icloud.com}{sideeqadeyemi@icloud.com}\\
\href{https://sadeeqgandalf.github.io}{sadeeqgandalf.github.io} \sep \href{https://github.com/sadeeqgandalf}{github.com/sadeeqgandalf} \sep
\href{https://www.linkedin.com/in/sadeeq-a-a90bb425a}{linkedin.com/in/sadeeq-a-a90bb425a}}

\section{Research Interests}
Multi-object tracking and segmentation (MOTS) through occlusion; 3D perception from cameras and lidar; state estimation and data
association; video foundation models (SAM 3); evaluation on public benchmarks.

\section{Education}
\entry{MS, Data Science and Engineering}{University of Maine, Orono, USA}{2025 -- present}
\entry{MS, Structural Engineering}{University of Maine, Orono, USA}{2023 -- 2025}
\entry{MSc, Project Management}{University of Portsmouth, United Kingdom}{2023}
\entry{BEng, Civil Engineering}{Federal University of Technology Akure, Nigeria}{2019}

\section{Research Experience}
\entry{Graduate Researcher, CVAR Lab (Computer Vision and Robotics)}{University of Maine}{2026 -- present}
How much identity can a tracker keep while an object is hidden? I designed and ran every experiment below.\par\vspace{1pt}
\begin{itemize}
  \item Benchmarked SAM 3 and SAM 3.1 with official HOTA: KITTI-MOTS val cars 82.6 and 74.6; MOTS20 train 72.6 and 61.5.
  \item Defined an occlusion measure from MOT17 visibility labels: SAM 3 kept identity through 94\% of 245 hides, SAM 3.1 63\%.
  \item Traced about half of SAM 3.1's deficit to its shipped tracking settings by reading the source code; a controlled run raised its car
        HOTA from 74.0 to 78.3, against SAM 3's 82.3. Isolated a separate fault that halts it on long, crowded sequences.
  \item Built a Kalman-filter identity layer on SAM's output: it repaired SAM 3.1 (identity through occlusion 63\% to 70\%, McNemar
        p = 0.004) but not SAM 3, which showed that motion outside the tracker says where an object is, not who it is.
  \item Mined 10,871 pedestrian occlusions in 1,256 NVIDIA driving clips with a 3D line-of-sight test. On that
        geometry an occlusion-aware tracker kept 57\% of identities through 1 to 5 second hides; a standard tracker kept 2\%.
  \item Built a simulation lab for controlled occlusion experiments (46 automated tests): identity kept rose from 0\% to 89\% as visibility
        reasoning, negative information and re-identification were added; MathWorks JPDA reached 57\%.
\end{itemize}

\entry{Robotics Research Assistant, CVAR Lab}{University of Maine}{2025 -- present}
\begin{itemize}
  \item Built mobile robots end to end, from sensor integration and wiring to ROS 2 nodes and lidar point-cloud perception.
  \item Implemented EKF localization fusing IMU, wheel odometry and lidar, with costmap planning and obstacle avoidance.
\end{itemize}

\section{Research Outputs}
\begin{itemize}
  \item Technical report: \emph{Solving occlusion in MOTS: SAM 3 and SAM 3.1 zero-shot}, CVAR Lab, October 2026. Available on request.
  \item Code: SAM 3 / SAM 3.1 benchmark, evaluation and filter layer, \href{https://github.com/sadeeqgandalf/av-stack}{github.com/sadeeqgandalf/av-stack}.
  \item Code and live lab: occlusion-aware tracking, \href{https://github.com/sadeeqgandalf/lunar-occlusion-tracking}{github.com/sadeeqgandalf/lunar-occlusion-tracking}.
\end{itemize}

\section{Industry Experience}
\entry{Materials Engineer}{R.W. Gillespie, Maine}{2024}
\entry{Machine Learning Engineer and Structural Steel Fabricator}{NewTech Steel, Nigeria}{2018 -- 2023}
\begin{itemize}
  \item Built and deployed a computer vision model for real-time fatigue and defect detection in steel, halving inspection time.
\end{itemize}
\entry{Civil Engineering Intern}{Federal Capital Development Authority, Abuja, Nigeria}{Feb -- Aug 2017}
\begin{itemize}
  \item Worked on the construction of wastewater treatment facilities for Nigeria's Federal Capital Territory.
\end{itemize}

\section{Technical Skills}
\skillrow{Methods}{multi-object tracking, Kalman filtering (EKF, UKF), data association (Hungarian, JPDA, MHT), 3D geometry and ray casting, re-identification}
\skillrow{Software and data}{Python, C++17, PyTorch, ROS 2, OpenCV, filterpy, MATLAB, TrackEval; KITTI, KITTI-MOTS, MOTS20, MOT17, NVIDIA PhysicalAI-AV}
\skillrow{Certifications}{AWS Certified Machine Learning Specialty; AWS Certified Solutions Architect, Associate}

\end{document}
